
\subsection{Đau đầu}

Đau đầu là triệu chứng của nhiều bệnh mà nguyên nhân rất phức tạp,
vì vậy không phải bất cứ đau đầu nào cũng xoa bóp được mà phải tuỳ
theo từng trường hợp cụ thể. Xoa bóp chỉ điều trị triệu chứng,
muốn điều trị nguyên nhân phải phối hợp các phương pháp khác của y
học hiện đại và y học cổ truyền. Muốn đem lại kết quả tốt cần phải
chỉ định cho đúng.

\subsubsection{Chỉ định}
\begin{itemize}
    \item Đau đầu do cảm mạo.
    \item Do thiểu năng tuần hoàn não.
    \item Do bệnh tâm căn suy nhược.
    \item Do bệnh tăng huyết áp mức độ nhẹ và trung bình.
    \item Do thoái hoá cột sống lưng, thắt lưng và cổ.
\end{itemize}

\subsubsection{Cách làm}
Bệnh nhân ngồi trên ghế tựa (nếu bệnh nhân khoẻ) hoặc bệnh nhân nằm
(nếu bệnh nhân yếu) lần lượt làm các thủ thuật sau:

\begin{longtblr}[
    caption = {Đau đầu},
]{
    width = \textwidth,
    colspec = {| X[1,c] |X[1, l]|X[2, l]|X[1, c]|X[7, l]|},
    hspan = minimal,
    rowhead = 1,
}
    \hline
    ID & Động tác & Vùng tác động & Số lần & Ghi chú
    \\ \hline
    1 & Chải & Đầu & 3-5 & {\
        Trước ra sau.\\
        Trên sang hai bên
    }
    \\ \hline
    2 & Miết & Trán & 3-5 & {\
        Miết từ ấn đường ra chân mày.
        Theo hình dẻ quạt ra hai bên.
    }
    \\ \hline
    3 & Miết & Huyệt & 3-5 & {\
        Ấn Đường, Thái Dương, Đầu Duy, Xuất Cốc, Phong Trì
    }
    \\ & & & & {\
        Ấn Đường
        \anduong
    }
    \\ & & & & {\
        Thái Dương
        \thaiduong
    } 
    \\ & & & & {\
        Đầu Duy
        \dauduy
    } 
    \\ & & & & {\
        Xuất Cốc
        \xuatcoc
    } 
    \\ & & & & {\
        Phong trì
        \phongtri
    }
    \\ \hline
    4 & Phân hợp & Trán & 3-5 & {\
        Phân hợp vùng trán, từ ấn đường ra chân mày, phân hợp hết trán, sử dụng ngón tay cái để tiến hành phân hợp.
    }
    \\ \hline
    5 & Véo & Véo từ đầu chân mày đến cuối chân mày & 3-5 & {\
        Dùng để xoá nhăn
    }
    \\ & & & & {\
        Ấn Đường
        \anduong
    }
    \\ & & & & {\
        Thái Dương
        \thaiduong
    }
    \\ \hline
    6 & Vỗ & Đầu, trán & 3-5 & {Vỗ vùng tiếp giáp đầu và trán}
    \\ \hline
    7 & Day & Hốc mắt & 3-5 & {\
        Dùng 3 đầu ngón tay giữa.
        Để vô lông mày, kéo ngược lên trán và day.
    }
    \\ \hline
    8 & Day & Đầu, trán & 3-5 & {\
        Dùng đầu ngón tay
    }
    \\ \hline
    9 & Chặt & Trán & 3-5 & {\
        Chặt xoè vùng đầu trán.
    }
    \\ \hline
    10 & Bóp & Đầu & 3-5 & {\
    }
    \\ \hline
    11 & Bấm/Day & Huyệt Ấn Đường & 1 (1 huyệt 1 phút) & {\
        Ấn Đường

        \anduong
    }
    \\ \hline
    12 & Bấm/Day & Huyệt Dương Bạch & 1 (1 huyệt 1 phút) & {\
        Dương Bạch

        \duongbach
    }
    \\ \hline
    12 & Bấm/Day & Huyệt Đầu Duy & 1 (1 huyệt 1 phút) & {\
        Đầu Duy

        \dauduy
    }
    \\ \hline
    13 & Bấm/Day & Huyệt Thái Dương & 1 (1 huyệt 1 phút) & {\
        Thái Dương

        \thaiduong
    }
    \\ \hline
    14 & Bấm/Day & Huyệt Hàm Yến & 1 (1 huyệt 1 phút) & {\
        Hàm Yên

        \hamyen
    }
    \\ \hline
    15 & Bấm/Day & Huyệt Bách Hội & 1 (1 huyệt 1 phút) & {\
        Bách Hội

        \bachhoi
    }
    \\ \hline
    16 & Bấm/Day & Huyệt Tứ Thần Thông & 1 (1 huyệt 1 phút) & {\
        Tứ Thần Thông

        \tuthanthong
    }
    \\ \hline
    17 & Bấm/Day & Huyệt Phong Phủ & 1 (1 huyệt 1 phút) & {\
        Phong Phủ

        \phongphu
    }
    \\ \hline
    18 & Bấm/Day & Huyệt Thiên Trụ & 1 (1 huyệt 1 phút) & {\
        Thiên Trụ

        \thientru
    }
    \\ \hline
    19 & Bấm/Day & Huyệt Phong Trì & 1 (1 huyệt 1 phút) & {\
        Phong Trì

        \phongtri
    }
    \\ \hline
    29 & Vận động & Cổ & 1 & {\
    }
    \\ \hline
\end{longtblr}


\subsubsection{Liệu trình}

Mỗi ngày xoa bóp một lần, mỗi lần 30 phút. Nếu do cảm mạo làm
1 đến 3 lần thì khỏi, còn các nguyên nhân khác chỉ điều trị triệu chứng,
vì vậy lúc nào đau thì xoa bóp.
